\section{Co-Occurrence of False Predictions} \label{sec:co-occurrence-of-false-predictions}
If different models both classify a portion of data as preictal, even though no seizure occurs afterward, this is interpreted as evidence for a proictal state, since preictal features were apparently present in the data.
To quantify the tendency of false predictions to co-occur across models, the \glsreset{cope}\gls{cope} \parencite{muller_coherent_2022} was calculated between the ensemble and \gls{cnn} scores for each patient.
Its definition follows.

Let $\vec{p}_i$ and $\vec{p}_j$ be the \textit{model's outputs} for models $i$ and $j$, respectively, and $\vec{L}$ the correct \textit{label} vector.\\
The \textit{prediction error} vector of model $i$ (analogously for model $j$) is defined as
\begin{equation}
    \vec{e}_i=\lvert\vec{p}_i-\vec{L}\rvert,
\end{equation}
the \textit{weight} for each sample $k$ as
\begin{equation} 
    w_k = \max(e_{i,k},e_{j,k}),\\
\end{equation}
the \textit{weighted mean} as
\begin{equation}
    \langle z\rangle_w = \frac{\sum_k w_k z_k}{\sum_k w_k},
\end{equation}
and the \textit{weighted standard deviation} as
\begin{equation}
    \sigma_w(z) = \sqrt{\langle z^2\rangle_w-\langle z\rangle_w^2}.
\end{equation}
The \textbf{\gls{cope}} is then
\begin{equation} \label{eq:cope}
    c_w=\frac{\left\langle\left(\vec{e}_i-\langle \vec{e}_i\rangle_w\right)\left(\vec{e}_j-\langle \vec{e}_j\rangle_w\right)\right\rangle_w}
    {\sigma_w(\vec{e}_i)\,\sigma_w(\vec{e}_j)}.
\end{equation}

This measure emphasizes coherence in false predictions by assigning higher weight to samples where at least one model has a large prediction error.
Consequently, samples with low errors contribute little, whereas samples with high errors dominate the estimate.